\clearpage\thispagestyle{empty}
\pdfbookmark[0]{Quellen zu Gauss an Bode (5. Mai 1812)}{bode1812may5-sources}
\begin{center}\textbf{Quellen zu Gauss an Bode (5. Mai 1812)}\end{center}
\noindent Diese KI-Übertragung ins heutige Deutsch folgt dem Brief von Gauß vom 5. Mai 1812 in den \href{https://gdz.sub.uni-goettingen.de/id/PPN236007467}{\textit{Werken}, Band~VI} (Göttingen, 1874), S.~351--354. Der Erstdruck erschien unter der Überschrift über Beobachtungen der Pallas, des zweiten Kometen von 1811 und seine Bahnelemente sowie Sternbedeckungen und Weiteres im \href{https://api.digitale-sammlungen.de/iiif/image/v2/bsb10538344_00005/full/1600,/0/default.jpg}{\textit{Astronomisches Jahrbuch für das Jahr 1815}}, herausgegeben von Johann Elert Bode (Berlin, 1812), S.~\href{https://api.digitale-sammlungen.de/iiif/image/v2/bsb10538344_00198/full/1600,/0/default.jpg}{190}, \href{https://api.digitale-sammlungen.de/iiif/image/v2/bsb10538344_00199/full/1600,/0/default.jpg}{191}, \href{https://api.digitale-sammlungen.de/iiif/image/v2/bsb10538344_00200/full/1600,/0/default.jpg}{192}, \href{https://api.digitale-sammlungen.de/iiif/image/v2/bsb10538344_00201/full/1600,/0/default.jpg}{193}, \href{https://api.digitale-sammlungen.de/iiif/image/v2/bsb10538344_00202/full/1600,/0/default.jpg}{194}. Das digitalisierte Exemplar stammt aus der Bayerischen Staatsbibliothek.
\par\smallskip\noindent Das Jahrbuch enthält auf S.~190 eine mit B. gezeichnete Anmerkung außerhalb des Textes der Werkausgabe. Auf S.~192 steht in der Perihellänge ein einfacher Strich nach 39. Auf S.~193 ist die Abweichung der Deklination vom 9. Dezember mit $+5''4$ angegeben. In der Werkausgabe stehen ein Doppelstrich nach 39 und $+3''4$. Die Übertragung folgt der Werkausgabe. KI-Übertragung und KI-Satz.
\clearpage
